\documentclass[11pt]{article}

% ====================================================
% PACKAGES
% ====================================================
\usepackage[margin=1in]{geometry}
\usepackage[T1]{fontenc}
\usepackage{lmodern}
\usepackage{microtype}
\usepackage{setspace}
\usepackage{parskip}

\usepackage{amsmath, amssymb, mathtools}
\usepackage{siunitx}
\usepackage{graphicx}
\usepackage{float}
\usepackage[justification=centering]{caption}
\usepackage{subcaption}

\usepackage{enumitem}
\usepackage{booktabs}

\usepackage{xcolor}
\usepackage{hyperref}

\usepackage{fancyhdr}
\usepackage{lastpage}

\usepackage{listings}
\usepackage{indentfirst}

% ====================================================
% Assignment Info
% ====================================================
\newcommand{\Name}{Arael Anaya}
\newcommand{\Assignment}{Test Rig Data Requirements}
\newcommand{\course}{Senior Design}

% ====================================================
% MATH SHORTHAND
% ====================================================
\newcommand{\st}{s\theta}
\newcommand{\ct}{c\theta}
\newcommand{\nh}[1]{\hat{n}_{#1}}
\newcommand{\dth}{\dot{\theta}}
\newcommand{\ddth}{\ddot{\theta}}
\newcommand{\Vg}{V_g}                 % gas volume -- distinct from spring PE, V
\newcommand{\DP}{\Delta P}
\newcommand{\Mk}{M_k}

% ====================================================
% PAGE STYLE
% ====================================================
\onehalfspacing%
\setlength{\parindent}{2em}

\pagestyle{fancy}
\fancyhf{}
\lhead{\Assignment\ |\ \course}
\rhead{\Name\ |\ \today}
\cfoot{Page \thepage\ of~\pageref{LastPage}}

\lstset{
    language=Matlab,
    basicstyle=\ttfamily\small,
    keywordstyle=\color{blue},
    commentstyle=\color{gray},
    stringstyle=\color{orange},
    numbers=left,
    numberstyle=\tiny,
    stepnumber=1,
    numbersep=5pt,
    frame=single,
    breaklines=true,
    tabsize=2,
    showstringspaces=false
}

\hypersetup{
    colorlinks,
    linkcolor={red!50!black},
    citecolor={blue!50!black},
    urlcolor={blue!80!black}
}

% ====================================================
% TITLE
% ====================================================
\title{Lagrangian Derivation for an $n$-Link Spring-Jointed Panel Array\\
\large with Inflatable Kink Actuation}
\author{\Name\\ \course}
\date{\today}

\begin{document}
\maketitle
\section{Purpose}

The scissor-boom model (\texttt{derivation/main.tex}, \texttt{src/}) leaves several inputs as
untested assumptions or placeholders. This document lists the data the test rig must collect to
replace those assumptions, validate the model, and size the stowed-envelope latch. The model treats
deployment as quasi-static: $N_s k\theta + 2N_k M_k(2\theta,\Delta P) = N_s\tau_f$.

\section{Model Inputs Needing Test Data}

\begin{center}
\begin{tabular}{llll}
\toprule
Quantity & Current value & Status & Section \\
\midrule
Hinge friction $F_{friction}$ & \SI{5}{N} & untested & \ref{sec:friction} \\
Spring rate $k$, preload & sized from friction & unverified & \ref{sec:spring} \\
Kink constant $a_{kink}$ & 1.0 & placeholder & \ref{sec:kink} \\
Closure form (C1 vs.\ C2) & assumed & unverified & \ref{sec:kink} \\
Tube radius $R_{tube}$ & \SI{0.1}{m} & placeholder & \ref{sec:kink} \\
Final angle $\theta_{min}$ & \SI{10}{\degree} & placeholder & \ref{sec:array} \\
Outward stow force & model prediction & unverified & \ref{sec:force} \\
\bottomrule
\end{tabular}
\end{center}

\section{Required Measurements}

\subsection{Hinge Friction (highest priority)}
\label{sec:friction}
\begin{itemize}[noitemsep]
    \item Breakaway (stiction) torque per hinge at stow ($\theta=\SI{90}{\degree}$).
    \item Running torque per hinge versus $\theta$, \SI{90}{\degree} to about \SI{10}{\degree}.
    \item Repeatability over multiple cycles, hinges, and single-hinge versus assembled.
    \item Dependence on hinge-pin load (panel mass, spring torque).
    \item Opening versus closing direction (model assumes pure Coulomb friction).
    \item Confirm effective moment arm, $\tau_f = F\,r_o$.
\end{itemize}

\subsection{Hinge Springs}
\label{sec:spring}
\begin{itemize}[noitemsep]
    \item Torque versus angle curve (model assumes $\tau=k\theta$).
    \item Spring rate $k$ (\si{N.m/rad}) and unit-to-unit variation.
    \item Preload or offset torque at $\theta=0$ (model assumes none).
    \item Peak stowed torque $k\theta_0$.
    \item Hysteresis or set after repeated cycles and long stow.
\end{itemize}

\subsection{Pressurized Kink Actuator}
\label{sec:kink}
\begin{itemize}[noitemsep]
    \item Kink moment $M_k$ versus fold angle $\beta=2\theta$ and $\Delta P$ (single-kink or blocked-moment test).
    \item Which closure fits: constant (C1, $M_k=a\Delta P R^3$) or vanishing fold (C2, $\propto\sin2\theta$).
    \item Empirical constant $a$, from several $\Delta P$ (and $R$ if possible).
    \item Moment near stow, to check C2's zero at \SI{90}{\degree}.
    \item Inflated tube radius $R$ versus $\Delta P$ and $\theta$.
    \item Optional: gas volume $V_g(\theta)$ to verify the volume-deficit model.
    \item Pressure drop and leakage (hold-pressure test).
    \item Fold location, wrinkling, and tube-to-link bond behavior.
\end{itemize}

\subsection{Whole-Array Deployment}
\label{sec:array}
\begin{itemize}[noitemsep]
    \item $\theta(t)$ from stow to $\theta_{min}$: spring-only (Phase 1) and spring plus pressure (Phase 2).
    \item Whether deployment completes, and the stall angle if not (compare with $\theta_1$).
    \item $\dot\theta$, $\ddot\theta$ from $\theta(t)$, to check the quasi-static assumption.
    \item Final angle reached versus $\theta_{min}=\SI{10}{\degree}$; effect of end stops.
    \item Pressure $P(t)$ logged alongside $\theta(t)$.
    \item Repeatability over multiple trials and re-stows; torque or force to re-stow.
\end{itemize}

\subsection{Outward Force and Latch Load}
\label{sec:force}
\begin{itemize}[noitemsep]
    \item Outward end force versus $\theta$ (model: $F=N k\theta/(nL\sin\theta)$, peak at stow).
    \item Stowed force at several lock points: far end, joint 1, link centre of mass.
    \item Peak force including dynamic and misalignment effects (scripts assume a $2\times$ placeholder).
\end{itemize}

\subsection{Geometry and Mass}
\begin{itemize}[noitemsep]
    \item Link mass (measured: \SI{40.4}{g}), panel mass, total array mass.
    \item As-built link length, hinge spacing, hole locations, number of links and springs.
    \item Centre of mass and, if desired, measured inertia of link and panel.
\end{itemize}

\subsection{Environment}
\begin{itemize}[noitemsep]
    \item Temperature (friction, spring rate, and membrane stiffness vary with it).
    \item Rig orientation and any gravity offloading (model is zero-g; bench is \SI{1}{g}).
    \item Shared time base across all channels.
\end{itemize}

\section{Test Methods}

\subsection{Common Rig Layout}

All tests use a flat base plate (aluminum or MDF) with the scissor array lying in the horizontal
plane, hinge axes vertical, so gravity acts normal to the motion and drops out of the torque balance.
The array's weight is carried on ball-transfer casters or PTFE pads. The caster/pad drag is measured
once (pull the array slowly with a force gauge, no springs fitted) and subtracted. One link end is
pinned to the plate at a fixed pivot; the far end node is free to travel along the $x$-axis.

Expected bench magnitudes (from \texttt{src/}, \texttt{testingParams.py}): $k\approx\SI{0.025}{N.m/rad}$,
peak stowed outward force $\approx\SI{0.15}{N}$ at the far end node, hinge friction torque
$\approx\SI{0.016}{N.m}$ (if $F=\SI{5}{N}$ at $r_o=\SI{3.2}{mm}$). These are small, so instruments must
resolve about \SI{0.01}{N} and \SI{0.001}{N.m}.

\subsection{Hinge Friction and Spring Torque (single-hinge fixture)}

\textbf{Setup.} One bench link is clamped flat to the base plate. The second link is attached through
the actual hinge pin, with the real washers and spring. The free link carries a hook at a marked
radius $r$ (for example \SI{200}{mm} from the pin axis). A digital force gauge (\SIrange{0}{5}{N},
\SI{0.01}{N} resolution) pulls on that hook with a string, kept perpendicular to the link. A
protractor scale, or a rotary encoder, sits under the pin to read $\theta$.

\textbf{Procedure.}
\begin{enumerate}[noitemsep]
    \item \emph{Friction only:} remove the spring. Pull slowly and steadily through \SIrange{90}{10}{\degree}.
          The peak before motion starts is breakaway; the steady reading is running friction.
          Torque is $\tau = F\,r$. Repeat in both directions, 10 cycles, on each hinge.
    \item \emph{Spring torque:} refit the spring. Load slowly from \SI{10}{\degree} to \SI{90}{\degree}
          and unload back, logging $F(\theta)$ at every \SI{5}{\degree}.
    \item The mean of the loading and unloading curves is the spring torque $\tau_s(\theta)$. Half the
          gap between them is the friction torque, a check on step 1. A straight-line fit gives $k$;
          the intercept at $\theta=0$ gives preload.
    \item \emph{Load dependence:} hang dead weights over a pulley, pulling the hinge radially along the
          link axis, and repeat step 1 to see how friction scales with pin load.
\end{enumerate}

\subsection{Kink Actuator (single-kink blocked-moment fixture)}

\textbf{Setup.} One kink: two short rigid arms with a pressurized tube segment bonded along both,
folded at the hinge. Arm A is bolted to the base plate. Arm B's tip presses against a miniature
compression or S-beam load cell ($\SIrange{0}{5}{N}$) at a known distance $d$ from the fold. The load
cell is bolted to a rigid post. The post sits on an arc of drilled holes (or a slotted arc) so the fold
angle $\beta$ can be set and locked at \SI{20}{\degree} steps from \SI{20}{\degree} to \SI{180}{\degree}.
The load cell is in line with the arm tip, with a spherical or rounded contact tip, so only the
normal force is read. Then $M_k = F\,d$.

The tube is fed from a regulator through a shut-off valve, with a low-range pressure transducer
(\SIrange{0}{7}{kPa} gauge, or \SI{1}{psi}) teed into the line at the tube inlet. The valve is closed after filling, so
$\Delta P$ is held constant while the load cell reads, and any leakage shows as a pressure decay.

\textbf{Procedure.} At each $\beta$, set $\Delta P$ to each of 4 or 5 levels, wait for the force
reading to settle, then log $F$, $\Delta P$ and tube width. Fit $M_k$ against $\beta$ at constant
$\Delta P$ to choose between the C1 and C2 closures. Fit $M_k$ against $\Delta P$ at constant $\beta$ to
get $a$ from $M_k = a\,\Delta P R^3$. Measure $R$ with calipers or a photo at each setting. Run the
fixture at several tube sizes if hardware allows, to confirm the $R^3$ scaling.

\textbf{Optional volume measurement.} Fill the tube with a calibrated syringe, hold $\Delta P$ constant,
and record displaced volume as $\beta$ changes; this gives $V_g(\theta)$ directly.

\subsection{Whole-Array Deployment}

\textbf{Setup.} The assembled 2-link array on the base plate, held stowed by the latch or a release pin.
An overhead camera (fixed, on a tripod above the plate) films the array with contrasting markers
(ArUco tags or dots) on each hinge and on the far end node. A rotary encoder or low-torque
potentiometer on one hinge, with a hollow bore through the $1/4$~in pin, is an optional cross-check;
since it adds friction, it is removed for the friction-sensitive runs. The tube pressure transducer
is logged by a DAQ (microcontroller or NI), and an LED flash visible to the camera marks a shared time
zero.

\textbf{Procedure.} Release from stow and film to full deployment or stall, then repeat at least 5
times. Track marker positions in post-processing to get $\theta(t)$, then smooth and differentiate to
get $\dot\theta$ and $\ddot\theta$. Compare stall angle, completion, and final angle with the model.
Run twice: springs only (tube unpressurized), then springs plus pressure. With the bench parameters
the model predicts a $\sim$\SI{2}{s} deployment, so \SI{60}{fps} or more is needed; the 1 hr and 11 hr
production timings are not reproduced on the bench, only the torque balance is validated.

\subsection{Outward Force and Latch Load}

\textbf{Setup.} A miniature compression load cell ($\SIrange{0}{5}{N}$, resolution
$\le\SI{0.01}{N}$) is bolted to a stop plate on a linear rail or slotted track along the $x$-axis of the
base plate. A rod with a rounded tip is threaded into the load cell and contacts the far end node, with
the load axis aligned with the deployment direction. The stop plate clamps at positions $x=nL\cos\theta$
for chosen angles. At each position the array is blocked by the load cell, which reads the outward force
directly. The load cell mount must be stiff, so the array can't move appreciably under load.

\textbf{Procedure.} For each $\theta$ from \SI{90}{\degree} (stowed, $x=0$) to \SI{10}{\degree}, approach
the position slowly from the stowed side and from the open side, reading $F$ each time. Friction acts
in opposite senses for the two approaches, so the pair brackets the frictionless value and gives the
upper bound the latch must hold. Repeat with the load cell at joint 1 and at a link centre of mass to
get force against lock location, and compare with $F=N k\theta/(c L\sin\theta)$.

\subsection{Geometry, Mass and Environment}

\begin{itemize}[noitemsep]
    \item Weigh each link, panel, and the assembled array on a scale (\SI{0.01}{g} resolution).
    \item Measure link length, hinge spacing and hole positions with calipers.
    \item Find centre of mass by balancing on a knife edge; optionally get inertia from a bifilar or
          torsion-pendulum swing test.
    \item Log room temperature with a thermocouple at the start of each test. If temperature
          sensitivity matters, repeat key tests in a warm or cold box.
\end{itemize}

\section{Instrumentation}

\begin{center}
\begin{tabular}{ll}
\toprule
Measurement & Suggested instrument \\
\midrule
Hinge angle & Rotary encoder or potentiometer; video tracking as cross-check \\
Hinge torque & Torque wrench or torque sensor \\
Outward / re-stow force & Load cell or force gauge \\
Tube pressure & Pressure transducer \\
Temperature & Thermocouple \\
Mass and geometry & Scale, calipers, tape \\
\bottomrule
\end{tabular}
\end{center}

\section{Test Sequence}

\begin{enumerate}[noitemsep]
    \item Single-hinge friction.
    \item Spring torque curve.
    \item Single-kink moment versus $\Delta P$.
    \item Spring-only array deployment (Phase 1).
    \item Spring plus pressure array deployment (Phase 2).
\end{enumerate}
Steps 1--3 replace placeholders in \texttt{src/testingParams.py} before the full-array tests.

\end{document}